\begin{figure}[ht]
\centering
\begin{tikzpicture}[x=1cm, y=1cm, font=\small]

% Residual phi(Z) = Z^3 - (1 - B) Z^2 + (A - 2B - 3B^2) Z - (AB - B^2 - B^3)
% for a near-critical case with three real roots: liquid, spurious, vapour.
% We use illustrative coefficients giving roots near Z = 0.07, 0.30, 0.62.
% Display: plotx = (Z - 0)*12 ; ploty = phi(Z)*60 (vertical units cm).

\def\sx{12.0}     % horizontal scale (Z in [0,0.75])
\def\sy{55.0}     % vertical scale for phi

% phi(Z) = (Z-0.07)(Z-0.30)(Z-0.62), expanded only for plotting.
% Axes
\draw[->, thick] (0, 0) -- ({0.78*\sx}, 0) node[right] {$Z$};
\draw[->, thick] (0, -1.7) -- (0, 2.1) node[above] {$\phi(Z)$};

% Curve
\draw[thick, engblue] plot[domain=0.02:0.74, samples=120]
  ({\x*\sx}, {(\x-0.07)*(\x-0.30)*(\x-0.62)*\sy});

% Roots
\fill[engred] ({0.07*\sx}, 0) circle (1.8pt);
\node[anchor=north, font=\scriptsize] at ({0.07*\sx}, -0.08) {$Z_{\ell}$};
\fill[enggray] ({0.30*\sx}, 0) circle (1.6pt);
\node[anchor=south, font=\scriptsize] at ({0.30*\sx}, 0.06) {$Z_{\mathrm{spur}}$};
\fill[engred] ({0.62*\sx}, 0) circle (1.8pt);
\node[anchor=north, font=\scriptsize] at ({0.62*\sx}, -0.08) {$Z_{v}$};

% A naive Newton step from a bad guess x0 = 0.45 (near the local max of phi
% between the spurious and vapour roots) overshoots far to the left.
\def\xa{0.45}
\pgfmathsetmacro{\pa}{(\xa-0.07)*(\xa-0.30)*(\xa-0.62)*\sy}
% slope of plotted curve at xa (numerically, central difference on display)
\pgfmathsetmacro{\hh}{0.001}
\pgfmathsetmacro{\pap}{((\xa+\hh-0.07)*(\xa+\hh-0.30)*(\xa+\hh-0.62)
  - (\xa-\hh-0.07)*(\xa-\hh-0.30)*(\xa-\hh-0.62))*\sy/(2*\hh*\sx)}
\fill[engamber] ({\xa*\sx}, \pa) circle (1.8pt);
\node[engamber, anchor=south west, font=\scriptsize] at ({\xa*\sx + 0.05}, {\pa+0.05}) {$Z^{(0)}$};
% tangent: zero crossing at plotx_zero
\pgfmathsetmacro{\xzero}{\xa*\sx - \pa/\pap}
\draw[thick, engamber] ({\xzero}, 0) -- ({\xa*\sx + 0.9}, {\pa + \pap*0.9});
\fill[engblack] ({\xzero}, 0) circle (1.4pt);
\node[engamber, anchor=south, font=\scriptsize] at ({\xzero}, 0.06) {full step};

\node[engblue, anchor=north east, font=\scriptsize] at ({0.74*\sx}, {(0.74-0.07)*(0.74-0.30)*(0.74-0.62)*\sy - 0.1}) {$\phi(Z)$};

\end{tikzpicture}
\caption[Cubic equation-of-state residual and the Newton trap]{The
cubic-equation-of-state residual $\phi(Z)$ for a near-critical
condition with three real roots: the liquid root $Z_{\ell}$, a
spurious middle root $Z_{\mathrm{spur}}$ with no physical meaning, and
the vapour root $Z_{v}$. A Newton step taken from a guess $Z^{(0)}$
near the interior local maximum, where $\phi'$ is small, overshoots far
past the nearest root, the divergent mode of section~5.4. The
globalised solver of the worked case damps the step and converges to
the physically correct root for the requested phase.}
\label{fig:vol02:ch05:eos-newton}
\end{figure}
